\section{Introduction}\label{sec:intro}

A phase-only spatial light modulator (SLM) placed in the pupil plane of a lens forms an array of optical traps whose number now reaches the thousands. The largest
published array holds about 6100 atoms in about 12,000 tweezers, and the same report states that the light field is one of the limits of
the machine: the trap depth varies by 11.4 percent (standard deviation) across the sites after weighted Gerchberg--Saxton optimization, while a small calibrated array reaches a relative
standard deviation of the intensity of 0.3 percent~\cite{manetsch,chew}. The spread therefore grows with the size of the array, and the light that reaches the traps is a budget that the
experiment spends on the uniformity of the sites. Programmable neutral-atom simulators of hundreds of atoms~\cite{ebadi} and the sorting of atoms into defect-free arrays make the quality of the hologram a question of the yield of the
experiment, not of its elegance.

The hologram is computed, not designed, and every method that computes it trades two quantities against each other. The efficiency $e$ is the fraction of the input power that lands in the traps. The uniformity $u$ measures how equal the traps are.
The foundational comparison of Di Leonardo et al.\ runs the random superposition, the Gerchberg--Saxton (GS) iteration, its adaptive additive variant, direct search and the weighted GS iteration on one ten by ten lattice and lists one
point of the plane $(e,u)$ for each: GS reaches an efficiency of 0.94 and a uniformity of 0.60 in 30 iterations, the adaptive additive variant 0.93 and 0.79, direct search 0.68 and 1.00 after $7.5\times10^{5}$ trial phases, and the weighted GS iteration 0.93 and 0.99~\cite{dileonardo}.
Direct search owes its position to a scalar weight $f$ that multiplies the spread in its cost, and the source reports it at two values of that weight~\cite{dileonardo,polin}. The phase-fixed form of weighted GS, which freezes the
phases of the traps after a few iterations, reaches a modulation efficiency of 91.2 percent and a non-uniformity below 0.5 percent on 1500 foci in simulation~\cite{kim2019}.

The gap that this paper addresses is the following. Each source reports one point, or at most two, of a trade-off curve, and it reports it for its own task, so a reader cannot say what an extra percent of uniformity costs in light on a given array, nor whether
a published point lies on the best attainable curve or inside it. The weights that appear in the cost of direct search and in the amplitude-freedom method of Pasienski and DeMarco~\cite{pasienski} are scalars fixed in advance. The
literature of multiobjective optimization has methods that compute a whole Pareto front of a smooth problem by solving a family of weighted problems, each started from the solution of its neighbour, and the measures that compare fronts, such as the hypervolume~\cite{zitzler2003,beume}.
No source of the holography literature read for this paper applies them to the trap-array hologram and reports a front for a given array against the points of the published methods run on the same array.

The proposed method computes that front. A weight $\lambda\in[0,1]$ combines two dimensionless terms into one loss: the squared deviation of the trap intensities from a flat target, and the negative of the light in the traps, both scaled to be of order one.
The problem is solved at $\lambda=0$, where the hologram is the refined output of phase-fixed weighted GS, and then at a sequence of increasing values of $\lambda$, each solved by conjugate-gradient iterations that start from the solution of the previous value. The gradient of the loss is closed form through the adjoint of the point-trap projection of the engine, so no automatic differentiation is needed. The set of points
$(e,u)$ of the solutions is the front, and its hypervolume summarizes it in one number that can be compared with the area covered by the published points.

The result on the ten by ten task of the engine is the following. Over \nSeeds{} seeds the front has a hypervolume of \hvFront{} (95 percent interval \hvFrontCi) against \hvPub{} (\hvPubCi) for the union of the points of the four published methods of the acceptance criterion, a difference of
\hvDiff{} (interval \hvDiffLo{} to \hvDiffHi) that is positive in \hvGeCount{} of \nSeeds{} seeds. The front dominates the points of Gerchberg--Saxton and of the adaptive additive iteration in every seed, but it dominates the point of weighted GS in only \domGSWSeeds{} and the point of phase-fixed WGS in only \domPFSeeds{} of \nSeeds{} seeds, and \frontDomMin{} to \frontDomMax{} of its \nLambda{} points are dominated by iterates of weighted GS run for 100 iterations. Against the benchmark of the studio, the stored front meets the published union on \benchPassThirty{} of \benchNThirty{} geometries at 30 iterations and on \benchPassHundred{} of \benchNHundred{} at 100 iterations, so the result is limited to the task on which it was measured. The contributions are the following.
\begin{enumerate}
\item A front-tracing method for point-trap holograms: a normalized weighted loss, its closed-form gradient, and a continuation over the weight with a warm start from the neighbouring value (Section~\ref{sec:method}).
\item An identity that rewrites the loss at every weight as the squared error against a uniformly raised target, which gives the weight a physical reading, turns the largest efficiency that is reached into the weight beyond which the target exceeds it (a restatement of a measurement, not a prediction), and suggests a schedule of weights (Proposition~\ref{prop:shift}).
\item Three further properties that are proved: the monotonicity of the light and of the error along the weights for global minimizers (Proposition~\ref{prop:mono}), two bounds that tie the measured spread and uniformity to the value of the loss (Proposition~\ref{prop:bound}), and the size of the starting gradient of a warm start (Proposition~\ref{prop:warm}).
\item A controlled comparison with the published methods run by the same engine on the same task and seeds, including the one-parameter fronts that those methods generate when their own parameter or their iteration count is varied (Sections~\ref{sec:design} and~\ref{sec:results}).
\item An ablation of twelve variants of the front, measured by a committed script, with the continuation, the schedule, the number of points, the steps per point, the start and the scaling of the loss each varied alone (Section~\ref{sec:ablation}).
\item The scope of the result, stated with the cases where it fails: the weighted GS and phase-fixed WGS points are dominated in a minority of seeds, the stored fronts of the benchmark meet the published union only part of the time, and the front is computed for the ideal device (Sections~\ref{sec:discussion} and~\ref{sec:limits}).
\end{enumerate}
Section~\ref{sec:related} places the method among the published ones, Section~\ref{sec:problem} fixes the model, the metrics and the multiobjective problem, Section~\ref{sec:method} defines the method, Section~\ref{sec:theory} gives the propositions that explain it, and
Sections~\ref{sec:design} to~\ref{sec:ablation} give the experiments. Section~\ref{sec:repro} lists the commands that rebuild every number and figure.
